\documentclass[CJK]{beamer}
\usepackage{CJKutf8}
\usepackage{beamerthemesplit}
\usetheme{Malmoe}
\useoutertheme[footline=authortitle]{miniframes}
\usepackage{amsmath}
\usepackage{amssymb}
\usepackage{graphicx}
\usepackage{color}
\usepackage{slashed}
\usepackage{simplewick}
\graphicspath{{../figures/}}
\def\be{\begin{equation}}
\def\ee{\nonumber\end{equation}}
\def\bea{\begin{eqnarray}}
\def\eea{\nonumber\end{eqnarray}}
\def\ii{{\dot{\imath}}}
\def\bch{\begin{CJK}{UTF8}{gbsn}}
\def\ech{\end{CJK}}
\def\bex{\begin{minipage}{0.3\textwidth}\includegraphics[width=1in]{jugelizi.png}\end{minipage}\begin{minipage}{0.6\textwidth}}
\def\eex{\end{minipage}}
\def\chtitle#1{\frametitle{\bch#1\ech}}
\def\skipline{{\vskip0.1in}}
\def\skiplines{{\vskip0.2in}}
\def\lagr{{\mathcal{L}}}
\def\hamil{{\mathcal{H}}}
\def\vecv{{\mathbf{v}}}
\def\vecx{{\mathbf{x}}}
\def\veck{{\mathbf{k}}}
\def\vecp{{\mathbf{p}}}
\def\vecn{{\mathbf{n}}}
\def\vecA{{\mathbf{A}}}
\def\vecP{{\mathbf{P}}}
\def\vecsigma{{\mathbf{\sigma}}}
\def\hatJn{{\hat{J_\vecn}}}
\def\hatJx{{\hat{J_x}}}
\def\hatJy{{\hat{J_y}}}
\def\hatJz{{\hat{J_z}}}
\def\hatj#1{\hat{J_{#1}}}
\def\hatphi{{\hat{\phi}}}
\def\hatq{{\hat{q}}}
\def\hatpi{{\hat{\pi}}}
\def\vel{\upsilon}
\def\Dint{{\mathcal{D}}}
\def\adag{{\hat{a}^\dagger}}
\def\bdag{{\hat{b}^\dagger}}
\def\cdag{{\hat{c}^\dagger}}
\def\ddag{{\hat{d}^\dagger}}
\def\hata{{\hat{a}}}
\def\hatb{{\hat{b}}}
\def\hatc{{\hat{c}}}
\def\hatd{{\hat{d}}}
\def\hatN{{\hat{N}}}
\def\hatH{{\hat{H}}}
\def\hatp{{\hat{p}}}
\def\Fup{{F^{\mu\nu}}}
\def\Fdown{{F_{\mu\nu}}}
\def\newl{\nonumber \\}
\def\SIkm{\mathrm{km}}
\def\SIyr{\mathrm{yr}}
\def\SIGyr{\mathrm{Gyr}}
\def\SIeV{\mathrm{eV}}
\def\SIGeV{\mathrm{GeV}}
\def\SIm{\mathrm{m}}
\def\SIcm{\mathrm{cm}}
\def\SIJ{\mathrm{J}}
\def\SIs{\mathrm{s}}
\def\SIkg{\mathrm{kg}}
\def\SIg{\mathrm{g}}
\def\vece{\mathrm{e}}
\def\bmat#1{\left(\begin{array}{#1}}
\def\emat{\end{array}\right)}
\def\bcase#1{\left\{\begin{array}{#1}}
\def\ecase{\end{array}\right.}
\def\calM{{\mathcal{M}}}
\def\calT{{\mathcal{T}}}
\def\calR{{\mathcal{R}}}
\def\barpsi{\bar{\psi}}
\def\baru{\bar{u}}
\def\barv{\bar{\upsilon}}
\def\bmini#1{\begin{minipage}{#1\textwidth}}
\def\emini{\end{minipage}}
\def\qeq{\stackrel{?}{=}}
\def\torder#1{\mathcal{T}\left(#1\right)}
\def\rorder#1{\mathcal{R}\left(#1\right)}


\title{Quantum Field Theory I \\ Lesson 05 - Gauge Field}
\author{}
\date{}


\begin{document}

\begin{frame}
 
\begin{center}
\begin{Large}
\bch
量子场论 I 

{\vskip 0.3in}

第五课 定域规范变换和规范场

\ech
\end{Large}
\end{center}

\vskip 0.2in

\bch
课件下载
\ech
https://github.com/zqhuang/SYSU\_QFTI

\end{frame}



\begin{frame}
\chtitle{定域规范变换(Local Gauge Transformation)}
\bch
上节课讲到，局域规范变换$\phi \rightarrow \phi e^{i\gamma(x)}$会改变不同时空点的$\phi$的相位差，从而改变作用量。联系不同时空点的$\phi$的算符为动量算符$-i\partial_\mu$，一个很自然的想法就是，能否把拉氏密度中的$-i\partial_\mu$算符替换成另一个矢量算符而使得作用量不变？

\skipline
答案是肯定的，通过定义“正则动量算符”：
$$-iD_\mu \equiv -i\partial_\mu + q A_\mu$$
并规定矢量场$A_\mu$在定域规范变换下按如下的规则进行变换：
$$ A_\mu \rightarrow A_\mu - \frac{1}{q} \partial_\mu \gamma\, ,$$
就可以满足作用量不变的要求。

\ech
\end{frame}

\begin{frame}
\chtitle{课堂互动}
\bch
按上述定义
$$D_\mu \equiv \partial_\mu + iq A_\mu$$
证明在定域规范变换$\phi \rightarrow e^{i\gamma(x)}\phi$, $A_\mu \rightarrow A_\mu - \frac{1}{q}\partial\gamma$下，拉氏密度
$$\langr = D_\mu\phi^\dagger D^\mu\phi - m^2\phi^\dagger \phi$$
是不变量。因此作用量也是不变量。

\skipline
$\phi \rightarrow e^{i\gamma(x)}\phi$在数学上称为一维幺正变换，简写为$U(1)$。与之对应的规范场$A_\mu$就叫做$U(1)$规范场。群论中U(1)属于阿贝尔群，所以$A_\mu$又称阿贝尔规范场。
\ech
\end{frame}

\begin{frame}
\chtitle{是否自寻烦恼？}
\bch
思考：为了使作用量形式上不变而引进的$A_\mu$场是“自寻烦恼”吗？
\skipline
关键在于：能否选取适当的局域规范变换使得$A_\mu$处处为零？
\ech
\end{frame}

\begin{frame}
\chtitle{课堂互动}
\bch
试证明：
\begin{itemize}
\item{如果存在$\gamma(x)$使得$A_\mu - \frac{1}{q}\partial_\mu\gamma = 0$处处成立，则
$F^{\mu\nu} \equiv \partial^\mu A^\nu - \partial^\nu A^\mu$ 处处为零。}
\item{若$F^{\mu\nu} \equiv \partial^\mu A^\nu - \partial^\nu A^\mu$处处为零，则存在$\gamma(x)$使得$A_\mu - \frac{1}{q}\partial_\mu\gamma = 0$处处成立}
\end{itemize}
\ech
\end{frame}

\begin{frame}
\chtitle{}
\bch
\ech
\end{frame}

\begin{frame}
\chtitle{}
\bch
\ech
\end{frame}

\begin{frame}
\chtitle{}
\bch
\ech
\end{frame}

\begin{frame}
\chtitle{}
\bch
\ech
\end{frame}

\begin{frame}
\chtitle{}
\bch
\ech
\end{frame}

\begin{frame}
\chtitle{}
\bch
\ech
\end{frame}

\begin{frame}
\chtitle{}
\bch
\ech
\end{frame}

\end{document}
